% ECONOMICS_PROBLEM: {"id": "D24-460", "title": "Misallocation versus measurement error", "jel_code": "D24", "source_id": "OP-0262"}
% FIRST_POSTED: 2026-10-09
% ATTEMPT_STATUS: unresolved
% ENTRY_KIND: research_attempt
% !TEX program = pdflatex
\documentclass[11pt]{article}
\usepackage[T1]{fontenc}
\usepackage[utf8]{inputenc}
\usepackage[margin=27mm]{geometry}
\usepackage{amsmath,amssymb,amsthm,mathtools}
\usepackage[colorlinks=true,linkcolor=blue,citecolor=blue,urlcolor=blue]{hyperref}
\allowdisplaybreaks
\newtheorem{theorem}{Theorem}
\newtheorem{proposition}[theorem]{Proposition}
\newtheorem{lemma}[theorem]{Lemma}
\newtheorem{corollary}[theorem]{Corollary}
\theoremstyle{remark}\newtheorem{remark}[theorem]{Remark}
\newcommand{\E}{\mathbb E}
\newcommand{\R}{\mathbb R}
\newcommand{\Ind}{\mathbf 1}
\newcommand{\Var}{\operatorname{Var}}
\newcommand{\Cov}{\operatorname{Cov}}
\newcommand{\TV}{\operatorname{TV}}
\newcommand{\KL}{\operatorname{KL}}
\newcommand{\supp}{\operatorname{supp}}
\DeclareMathOperator*{\argmin}{arg\,min}
\author{GPT 6 Astra Ultra}
\date{9 October 2026}
\title{Auditing the counterfactual gain from reallocation\\\large A CES benchmark, exact error sets, and failure under persistent mismeasurement}
\begin{document}
\maketitle
\noindent\textbf{Problem ID:} \texttt{D24-460}. \textbf{Status:} Unresolved; rigorous results in explicitly declared models.
\section{The counterfactual being measured}
Dispersed measured revenue products do not alone identify a productivity gain. We derive the gain for the catalogue's CES quantity aggregator and common Cobb--Douglas exponents, under an explicit restricted benchmark with no adjustment or transport costs. We then show how independent audits identify this gain and how unrestricted persistent errors prevent identification even with arbitrarily long panels.

Bils, Klenow and Ruane \cite{bkr} were consulted in their original February 2020 Census working paper, abstract and Sections 1--2. Their introduction distinguishes true and measured marginal products, and states restrictions needed for their growth-based correction. Our repeated-log-audit model differs from their correction and is not claimed as their result. It is a tractable instance of the editorial question.

There are $N$ plants, known $\alpha,\gamma,\eta>0$ with $\alpha+\gamma+\eta=1$, and true outputs
\[
 q_i=A_i k_i^\alpha\ell_i^\gamma m_i^\eta,
 \quad A_i,k_i,\ell_i,m_i>0.
\]
The known fixed aggregate resources are $(K,L,M)$, equal to sums of the current inputs. Let $\epsilon>1$, $\rho=1-1/\epsilon\in(0,1)$ and $\omega_i>0$. The same final-good index is used before and after reallocation:
$Y=(\sum_i\omega_iq_i^\rho)^{1/\rho}$.
The benchmark removes all constraints causing cross-plant factor misallocation, holding these resources and technologies fixed. Common markups do not create relative factor wedges here; heterogeneous markups, input-supply constraints, or unavoidable costs would require a different feasible benchmark. The theorem is about physical quantities, not a revenue-only production-function identification claim.

\section{An exact efficient-output formula}
Define
\[
 H(A)=\left(\sum_i\omega_i^\epsilon A_i^{\epsilon-1}\right)^{1/(\epsilon-1)},
 \qquad B=K^\alpha L^\gamma M^\eta.
\]
\begin{theorem}[Efficient allocation]\label{efficient}
The maximum of the same final-good index over all nonnegative allocations with the specified aggregate resources is
\[
 Y^{\rm eff}=B H(A).
\]
It is attained by assigning the same fraction of every input to plant $i$:
\[
 (k_i^*,\ell_i^*,m_i^*)=s_i(K,L,M),\qquad
 s_i=\frac{\omega_i^\epsilon A_i^{\epsilon-1}}
              {\sum_j\omega_j^\epsilon A_j^{\epsilon-1}}.
\]
Thus the gain is $G=BH(A)/Y-1\geq0$ for every feasible current allocation.
\end{theorem}
\begin{proof}
Write normalized input shares $u_i=k_i/K$, $v_i=\ell_i/L$, $t_i=m_i/M$ and $z_i=u_i^\alpha v_i^\gamma t_i^\eta$. Generalized H\"older gives
$\sum_i z_i\leq(\sum_i u_i)^\alpha(\sum_i v_i)^\gamma(\sum_i t_i)^\eta=1$.
For completeness, its finite nonnegative version follows from weighted Young's inequality $r^\alpha s^\gamma t^\eta\leq\alpha r+\gamma s+\eta t$ after normalizing each sequence's relevant power sum; Young follows from concavity of $\log$ (with zero coordinates by a limit).
Therefore
\[
 Y^\rho=B^\rho\sum_i c_i z_i^\rho,
                    \qquad c_i=\omega_i A_i^\rho.
\]
The strictly concave function $\sum_i c_i z_i^\rho$ on $\{z\geq0:\sum_i z_i\leq1\}$ is increasing, so the optimum has total one. Its maximizer is interior since a zero component has infinite right derivative. The first-order equations $\rho c_i z_i^{\rho-1}=\lambda$ give
$z_i=c_i^{1/(1-\rho)}/\sum_jc_j^{1/(1-\rho)}=s_i$.
Concavity verifies global optimality. Substitution gives
$\max\sum_i c_i z_i^\rho=(\sum_i c_i^{1/(1-\rho)})^{1-\rho}$;
raising to $1/\rho$ gives $H(A)$. The common input shares $s_i$ realize $z_i=s_i$, so the bound is feasible. The current allocation is in the same feasible set, giving $G\geq0$.
\end{proof}

\section{What repeated audits identify}
For each of the four true log quantities $z_{ih}\in\{\log q_i,\log k_i,\log\ell_i,\log m_i\}$, suppose $R$ contemporaneous measurements satisfy
\[
 V_{ihr}=z_{ih}+e_{ihr},\quad r=1,\ldots,R.
\]
The true plant quantities are fixed during these audits. Assume errors are independent over $r$, have mean zero and are sub-Gaussian with supplied scale $\sigma_h$:
$\E e^{t e_{ihr}}\leq e^{t^2\sigma_h^2/2}$ for every real $t$.
Errors may be correlated across plants and quantity types within an audit round. Let $\widehat z_{ih}=R^{-1}\sum_rV_{ihr}$, exponentiate to estimate quantities, and put
$\widehat A_i=\widehat q_i/(\widehat k_i^\alpha\widehat\ell_i^\gamma\widehat m_i^\eta)$.
Use the known resource totals to define the raw ratio
$\widehat R_G=BH(\widehat A)/Y(\widehat q)>0$ and
$\widehat G=\max(0,\widehat R_G-1)$.

\begin{lemma}[Logarithmic stability]\label{stable}
If $\max_i|\widehat z_{ih}-z_{ih}|\leq d_h$ for all four types, then
\[
 |\log(1+\widehat G)-\log(1+G)|
 \leq D:=2d_q+\alpha d_k+\gamma d_\ell+\eta d_m.
\]
\end{lemma}
\begin{proof}
The log technology errors are bounded by
$d_A=d_q+\alpha d_k+\gamma d_\ell+\eta d_m$.
Monotonicity and degree-one homogeneity of $H$ give
$e^{-d_A}H(A)\leq H(\widehat A)\leq e^{d_A}H(A)$.
The same properties of the CES output index give
$e^{-d_q}Y(q)\leq Y(\widehat q)\leq e^{d_q}Y(q)$.
Taking logarithms of their ratio bounds its error by $d_A+d_q=D$. Since $\log(1+G)\geq0$ and $\log(1+\widehat G)=\max(0,\log\widehat R_G)$, projection onto $[0,\infty)$ cannot increase its error.
\end{proof}

\begin{theorem}[Audit confidence and consistency]\label{audit}
For any $\zeta\in(0,1)$, take
\[d_h=\sigma_h\sqrt{2\log(8N/\zeta)/R},\]
and define $D$ as in Lemma~\ref{stable}. With probability at least $1-\zeta$, the true gain belongs to
\[
 \left[\max\{0,(1+\widehat G)e^{-D}-1\},                     (1+\widehat G)e^D-1\right].
\]
For fixed $N$ and the declared technology and resource parameters, $\widehat G\to G$ in probability as $R\to\infty$. Correlation within an audit round does not invalidate these conclusions.
\end{theorem}
\begin{proof}
Independence across rounds and the moment-generating-function bound imply
$\Pr(|\widehat z_{ih}-z_{ih}|>d_h)\leq2\exp[-Rd_h^2/(2\sigma_h^2)]=\zeta/(4N)$, with the zero-scale case interpreted as an exact measurement. Union bounding over $4N$ coordinates yields the event in Lemma~\ref{stable}. Exponentiating its conclusion gives the interval. To obtain convergence, each fixed coordinate sample average converges in probability by its variance bound, so the logarithmic error tends to zero. Continuity of the exponential at the finite true gain proves the claim.
\end{proof}

This uncertainty concerns the gain, not merely dispersion in a proxy. A confidence statement conditional on wrong CES weights or production elasticities is not robust to that misspecification.

\section{Sharp bounds from bounded audit errors}
Suppose instead that a validation protocol supplies deterministic log-error intervals $z_{ih}\in[\underline z_{ih},\overline z_{ih}]$. Intersect their product with
\[
 \sum_i e^{z_{ik}}=K,\qquad
 \sum_i e^{z_{i\ell}}=L,\qquad
 \sum_i e^{z_{im}}=M.
\]
Call the resulting set $\mathcal Z$. It is compact, though generally nonconvex. Output and inputs are strictly positive on it. For each feasible $z$, derive $A_i$ from the production identity and evaluate $G(z)$ using Theorem~\ref{efficient}.

\begin{proposition}[An exact finite-dimensional gain set]
If the only restrictions on true quantities and measurement errors are the preceding bounds, resource identities and production equations, the sharp gain set is
$\{G(z):z\in\mathcal Z\}$. When $\mathcal Z\neq\varnothing$, its minimum and maximum are attained. They bound every compatible gain, but the image need not be asserted to be an interval without a connectedness argument. If the audit box jointly covers the true log vector with probability $1-\zeta$, the two optimized endpoints cover the true gain with at least that probability.
\end{proposition}
\begin{proof}
Every compatible model has a true vector in $\mathcal Z$, giving one inclusion. Conversely, every $z\in\mathcal Z$ defines positive true quantities and technologies satisfying production and resources. Setting each measurement error equal to the observed log quantity minus its chosen true log value realizes the observed records within the stated bounds. There are no further error restrictions in this proposition, so every retained value is attainable. Continuity of $G$ and compactness give extrema. On the audit coverage event the true vector is feasible, which proves coverage.
\end{proof}

\section{Two ways repeated records can mislead}
\begin{proposition}[A constant panel does not identify misallocation]
With two plants, $\omega_1=\omega_2=1$, $\rho=1/2$, $A_1=A_2=1$ and $K=L=M=1$, identical measured quantities $\widetilde q_i=\widetilde k_i=\widetilde\ell_i=\widetilde m_i=1/2$ in every period are compatible both with $G=0$ and with gains arbitrarily close to one, if persistent correlated log errors are unrestricted.
\end{proposition}
\begin{proof}
For any $t\in(0,1)$ assign every input share $(t,1-t)$. Constant returns give $q=(t,1-t)$. True output is
$Y=(\sqrt t+\sqrt{1-t})^2$ and efficient output is $2$. Set every type of measurement error for plant $i$ in every period to $\log[(1/2)/q_i]$. All measured records are then identically one half. At $t=1/2$ the gain is zero; as $t\downarrow0$, output tends to one and the gain tends to one. The same measurement law is observed for every panel length.
\end{proof}

Independent duplicate measurements can identify a variance without identifying every firm's realization. If random plants satisfy $V^{(1)}=Z+e_1$, $V^{(2)}=Z+e_2$, with finite second moments, $\E[e_r\mid Z]=0$ and $\E[e_1e_2\mid Z]=0$, then
$\Cov(V^{(1)},V^{(2)})=\Var(Z)$ by expanding the covariance and conditioning on $Z$. If both measurements instead include a shared persistent component $B$, the data depend on $Z+B$ and cannot generally distinguish that component from truth. Even identification of $\Var(Z)$ does not determine the nonlinear sum of exponentials in $H(A)$, so a variance correction is not by itself a gain identification theorem.

The remaining economic work is to justify independent validation, distinguish physical productivity from demand, and retain unavoidable costs in the feasible benchmark. Eliminating an unknown subset of wedges, or allowing endogenous resources and heterogeneous markups, changes the counterfactual itself. The restricted formulas give exact targets and audit requirements for that further work, not a universal industry-level gain.

\section{Completion Estimate}
65\%

This is a post hoc, heuristic estimate for the full catalogue target.
The CES benchmark yields an exact efficient allocation and productivity gain, with repeated audit confidence bounds and sharp bounded error gain sets. Independent validation, demand heterogeneity, persistent correlated measurement error, heterogeneous markups, and the counterfactual retaining unavoidable adjustment and transport costs remain unsettled.

\begin{thebibliography}{9}
\bibitem{bkr} M. Bils, P. J. Klenow and C. Ruane (2020), \emph{Misallocation or Mismeasurement?}, Census CES 20-07, February 2020. Original PDF, abstract and Sections 1--2 inspected 9 October 2026. \url{https://www2.census.gov/ces/wp/2020/CES-WP-20-07.pdf}.
\end{thebibliography}
\end{document}
